\section{Part V --- The modifications, one by one}

\noindent\textbf{Principle.} The original files are left untouched except for four small \emph{registrations} (M12). Every new behaviour lives in new files whose names end in \code{\_net} or \code{network}, so that \code{git diff} shows exactly what was added and the tree sampler keeps working unchanged (\code{ENVIRONMENT\_TYPE: ONE\_STEP\_BINARY\_TREE}). Table~\ref{tab:mods} is the map.

\begin{table}[h]\centering\small
\begin{tabular}{lp{7.2cm}p{5.6cm}}\toprule
\textbf{Id} & \textbf{File (new unless marked)} & \textbf{Content} \\ \midrule
M1 & \pfile{src/env/phylo_network_env.py} & node / lineage / state / network data structures \\
M2 & \pfile{src/env/phylo_network_env.py} & initial state and the two-version feature tensor \\
M3 & \pfile{src/env/phylo_network_env.py} & level-1 masks \\
M4 & \pfile{src/env/phylo_network_env.py} & transitions \textsc{Merge} / \textsc{Reticulate}, ordering keys \\
M5 & \pfile{src/env/phylo_network_env.py} & feature updates = exact likelihood, reward \\
M6 & \pfile{src/env/phylo_network_env.py} & number of parents, backward sampling from a network \\
M7 & \pfile{src/env/phylo_network_env.py} & policy inputs: normalisation, flags, masks, padding \\
M8 & \pfile{src/model/network_model/one_step_network_model.py}, \pfile{src/model/edges_model/continuous/network_heads.py} & policy with two action types; heads for lengths and $\theta$ \\
M9 & \pfile{src/utils/level1_counts.py} & exact class sizes; topology prior normaliser \\
M10 & \pfile{src/gfn/tb_gfn_net.py} & generator: routing of actions, TB loss \\
M11 & \pfile{src/gfn/rollout_worker_net.py}, \pfile{src/gfn/training_data_loader_net.py}, \pfile{src/gfn/gfn_evaluator_net.py} & variable-length rollouts, replay buffer, evaluation \\
M12 & \pfile{src/configs/defaults.py}, \pfile{src/env/__init__.py}, \pfile{src/gfn/build.py}, \pfile{src/utils/utils.py} (modified); \pfile{train_net.py}, \pfile{src/configs/network_cfgs/*.yaml}, \pfile{tools/*}, \pfile{tests/*} & registrations, config, training script, data simulation, tests \\
\bottomrule
\end{tabular}
\caption{Where each modification lives.}\label{tab:mods}
\end{table}

% ---------------------------------------------------------------------------
\subsection{M1 --- Data structures: nodes with two parents}
\textbf{What.} PhyloGFN stores trees as \code{ete3.TreeNode}, which has exactly one parent per node. A reticulation node needs two. \code{NetNode} is a minimal DAG node; \code{Lineage} is an entry of the state list (a real root, or a bare slot); \code{NetworkState} replaces \code{PhylogenticTreeState}; \code{PhyloNetwork} replaces \code{PhylogeneticTree} (same interface: \code{log\_score}, \code{signature}, ordering for the replay heap).
\begin{lstlisting}
class NetNode(object):
    """kind: 'leaf' | 'tree' | 'ret'
       ret: one child (child_dist), two parents = two SLOTS parents[k], parent_dists[k];
            theta = inheritance weight of slot 0, 1 - theta of slot 1."""
    __slots__ = ('idx', 'kind', 'children', 'child_dists', 'parents', 'parent_dists',
                 'theta', 'child_dist', 'tokens', 'ret_id')

class Lineage(object):
    """node: NetNode; slot: None (real lineage) or 0/1 (bare parent slot of node = h);
       open_h: the unclosed reticulation carried (None if closed);
       tokens: ownership tokens; key = min(tokens) (unique -> sorted state list)."""

class NetworkState(object):
    def __init__(self, lineages, next_idx, n_ret, log_reward=None, log_score=None):
        self.lineages = lineages; self.next_idx = next_idx; self.n_ret = n_ret
        self.is_done = (len(lineages) == 1 and lineages[0].open_h is None)
\end{lstlisting}
\textbf{Why.} Everything downstream (transitions, likelihood, backward sampling, canonical signatures) needs to know, for a reticulation, \emph{which} parent is which and what $\theta$ belongs to which slot. \code{network\_canonical\_string} writes a reticulation subtree as \code{\#<subtree>} at both parent positions; on level-1 networks with labelled leaves this string is a faithful topology identifier (used by the replay buffer instead of ete3's \code{get\_topology\_id}).

% ---------------------------------------------------------------------------
\subsection{M2 --- Initial state and the feature tensor}
\textbf{What.} \code{get\_initial\_state} creates $N$ leaf lineages with token sets $\{0\},\dots,\{N-1\}$. \code{initial\_features(B)} returns a tensor of shape $[B,\,l_{\max},\,2,\,m,\,4]$ (two versions per lineage) filled with ones, with the one-hot sequences copied into \emph{both} versions of the first $N$ rows.
\begin{lstlisting}
self.l_max = self.n_leaves + self.r_max            # a RETICULATE adds one lineage
rows, cols = torch.triu_indices(self.l_max, self.l_max, offset=1)
self.n_pairs = len(rows);  self.n_actions = self.n_pairs + self.l_max   # [pairs | singles]

def initial_features(self, batch):
    feats = torch.ones(batch, self.l_max, 2, self.m, self.c, dtype=..., device=...)
    feats[:, :self.n_leaves, 0] = self.seq_arrays
    feats[:, :self.n_leaves, 1] = self.seq_arrays
    return feats
\end{lstlisting}
\textbf{Why.} PhyloGFN's tensor is $[B,l,m,4]$ with $l$ shrinking by one every step for \emph{all} trajectories. Ours must (i) hold two versions and (ii) keep a fixed width $l_{\max}=N+R_{\max}$ because the number of lineages differs between trajectories of the same batch; unused rows are ones and are hidden from the policy by a padding mask (M7). The action space is laid out once, over $l_{\max}$: pair $k\mapsto(\code{rows}[k],\code{cols}[k])$, single $n_{\text{pairs}}+i\mapsto$ lineage $i$.

% ---------------------------------------------------------------------------
\subsection{M3 --- Level-1 masks}
\textbf{What.} \code{legal\_actions(state)} returns a boolean vector over the padded action space; \code{merge\_legal(u,v)} is the pair rule.
\begin{lstlisting}
def legal_actions(self, state):
    L = state.lineages; n = len(L)
    mask = np.zeros(self.n_actions, dtype=bool)
    if state.is_done: return mask
    for i in range(n):
        for j in range(i + 1, n):
            mask[self.pair_index[(i, j)]] = self.merge_legal(L[i], L[j])
    n_free = sum(1 for l in L if l.open_h is None)
    if state.n_ret < self.r_max:
        for i in range(n):
            if L[i].open_h is None and n_free >= 2:     # x free AND another free lineage exists
                mask[self.n_pairs + i] = True
    return mask

@staticmethod
def merge_legal(u, v):
    if u.open_h is not None and v.open_h is not None:
        if u.open_h is not v.open_h: return False     # two different open cycles: level-2
        if u.is_bare and v.is_bare:  return False     # the two bare slots of one h: parallel edges
    return True
\end{lstlisting}
\textbf{Why.} See Part~IV.2. The three rules are $O(1)$ per candidate action because $\mathrm{open}(u)$ is stored on the lineage, never recomputed from the graph. In PhyloGFN every pair is legal, so the only mask is the padding mask; here the mask is also the definition of the class.

% ---------------------------------------------------------------------------
\subsection{M4 --- Transitions and ordering keys}
\textbf{What.} \code{apply\_action(state, action, feats)} implements both actions on the Python state and (optionally) on the feature tensor; it returns \code{(new\_state, new\_feats, log\_score)} with \code{log\_score} set only at a terminal state. The \textsc{Merge} branch is PhyloGFN's \code{update\_state} plus the slot bookkeeping; the \textsc{Reticulate} branch is new.
\begin{lstlisting}
# MERGE(i, j): new tree node w; a bare slot attaches the reticulation node itself
for lin, d, slot_side in ((u, bl, 0), (v, br, 1)):
    child = lin.node
    w.children.append(child); w.child_dists.append(d)
    if lin.is_bare:                                # fill the slot: edge w -> h of length d
        child.parents[lin.slot] = w; child.parent_dists[lin.slot] = d
open_h = None if case in (0, 3) else (u.open_h if case == 1 else v.open_h)
new_lin = Lineage(w, open_h=open_h, tokens=u.tokens | v.tokens)
new_L = L[:i] + [new_lin] + L[i + 1:j] + L[j + 1:]      # result at position i, j removed

# RETICULATE(i): h above lineage x, two bare slots p1 (position i) and p2 (appended)
h = NetNode(state.next_idx, 'ret'); h.children = [x.node]; h.child_dist = b_hx; h.theta = theta
fresh = self.n_leaves + state.n_ret                     # ordering token of slot 1
p1 = Lineage(h, slot=0, open_h=h, tokens=x.tokens)
p2 = Lineage(h, slot=1, open_h=h, tokens=[fresh])
new_L = L[:i] + [p1] + L[i + 1:] + [p2]
\end{lstlisting}
\textbf{Why.} \code{case} (0 both closed, 1 only $u$ open, 2 only $v$ open, 3 closing) decides both the open status of the result and the feature rule (M5). The token discipline replaces \code{min\_seq\_idx}: keys are unique and the list stays sorted, which is what the backward sampler needs to reconstruct the same action indices (Part~IV.4).

% ---------------------------------------------------------------------------
\subsection{M5 --- Feature updates: the exact likelihood and the reward}
\textbf{What.} Three private methods replace \code{compute\_partial\_prob}'s role for networks.
\begin{lstlisting}
def _merge_features(self, feats, i, j, case, bl, br, last_step):
    uA, uB = feats[i, 0], feats[i, 1];  vA, vB = feats[j, 0], feats[j, 1]
    M = self.evolution_model.get_transition_matrices(torch.tensor([bl, br]))
    if last_step:      # PhyloGFN root convention: one prior on b_l + b_r (pulley principle)
        tu = lambda f: f @ M[0]; tv = lambda f: f @ M[1]; prior = exp(log_prior(bl + br) / m)
    else:              # edge priors folded into the features, site by site (paper eq. 4)
        pl = exp(log_prior([bl, br]) / m)
        tu = lambda f: (f @ M[0]) * pl[0]; tv = lambda f: (f @ M[1]) * pl[1]; prior = 1.0
    if case == 0:   wA = tu(uA) * tv(vA) * prior;                 wB = wA
    elif case == 1: wA = tu(uA) * tv(vA) * prior;                 wB = tu(uB) * tv(vA) * prior
    elif case == 2: wA = tu(uA) * tv(vA) * prior;                 wB = tu(uA) * tv(vB) * prior
    else:           wA = (tu(uA) * tv(vB) + tu(uB) * tv(vA)) * prior;  wB = wA   # closing merge
    ... row j removed, row i <- (wA, wB), a row of ones appended to keep l_max rows

def _reticulate_features(self, feats, i, n, theta, b_hx):
    Lh = self._transport(feats[i, 0][None], torch.tensor([b_hx]))[0]   # M(b_hx) L_x * prior^(1/m)
    new[i, 0] = theta * Lh;        new[i, 1] = ones          # slot 0 at position i
    new[n, 0] = (1 - theta) * Lh;  new[n, 1] = ones          # slot 1 appended at position n

def log_topology_prior_unnormalised(self, n_ret):
    if self.prior_type == 'UNIFORM_R':                                   # p(R) ~ exp(-λR), uniform within R
        return -self.lambda_r * n_ret - self.log_class_size[n_ret]
    return -self.lambda_r * n_ret                                        # EXP_R: p(G) ~ exp(-λR) over the class

def _terminal_log_score(self, root_feat, n_ret):
    log_p = torch.sum(torch.log(torch.sum(root_feat / 4.0, -1)), -1)   # sum_i log sum_a A^i(a)/4
    return float(log_p) + self.log_topology_prior_unnormalised(n_ret)   # + log p(G)
\end{lstlisting}
\textbf{Why.} This is Proposition~\ref{prop:mix} in code. Note what did \emph{not} change: the transition matrices, the edge prior, the per-site folding of $P(b)^{1/m}$, the root formula $\sum_i\log\sum_a L^i(a)/4$ and the last-step convention are all PhyloGFN's --- with $R_{\max}=0$ the code computes exactly PhyloGFN's reward (Part~VI.4). The reward reshaping \code{(C + log\_score)/scale} is reused from \code{PhyloTreeReward}.

% ---------------------------------------------------------------------------
\subsection{M6 --- Backward policy: counting parents, sampling backward}
\textbf{What.} \code{NetworkState.num\_parents()} and \code{PhyloNetworkEnv.sample\_backward\_from\_network(net)}.
\begin{lstlisting}
def num_parents(self):
    n = 0; seen_ret = {}; n_free = 0
    for l in self.lineages:
        if l.is_bare: seen_ret[l.node.idx] = seen_ret.get(l.node.idx, 0) + 1
        else:
            if l.open_h is None: n_free += 1
            if l.node.kind == 'tree': n += 1                   # un-merge
    if n_free >= 1:
        n += sum(1 for c in seen_ret.values() if c == 2)       # un-reticulate (only if legal forward)
    return n
\end{lstlisting}
\code{sample\_backward\_from\_network} walks backward from the finished network: at every step it lists the undo moves (un-merge any tree-root lineage; un-reticulate any $h$ whose two slots are both bare, provided a free lineage exists), records their number ($=|\mathrm{Pa}|$), picks one uniformly, and stops at the leaves. It then \emph{replays forward}, matching the original nodes to freshly created ones (\code{new\_of}) to recover the pair / lineage indices under the sorted convention, and attaches the stored lengths and $\theta$'s. It returns \code{(actions, log\_paths\_pb)} exactly like \code{sample\_backward\_from\_tree}.\\
\textbf{Why.} Uniform $\PB$ requires the true number of parents in the state graph (Part~IV.5); the replay-buffer training and the Pearson-$r$ evaluation both need backward-sampled trajectories of a \emph{given} network. \code{test\_backward\_sampling\_roundtrip} checks that the replay yields the same canonical network, the same score and the same $|\mathrm{Pa}|$ sequence as the forward code.

% ---------------------------------------------------------------------------
\subsection{M7 --- What the policy sees}
\textbf{What.} \code{prepare\_rollout\_inputs(feats, states, random\_spec, input\_actions)} normalises both versions per site (as PhyloGFN's \code{NORMALIZE\_LIKELIHOOD}), flattens them, appends four flags per lineage, and returns the padding mask, the action mask and the pair index tables.
\begin{lstlisting}
x = feats / feats.sum(-1, keepdim=True); x[isnan] = 1/4          # per-site normalisation of A and B
x = x.reshape(b, self.l_max, -1)                                  # [b, l_max, 2*m*4]
flags[k, i] = [has_open, is_bare, side (slot 0 / slot 1), share (theta or 1 - theta)]
x = torch.cat([x, flags], dim=-1)
pad = torch.arange(self.l_max)[None] >= nb[:, None]               # True = padding
masks[k] = self.legal_actions(s)                                  # True = legal
\end{lstlisting}
\textbf{Why.} The Transformer needs a set of fixed-size vectors; padding rows are excluded by \code{key\_padding\_mask}, illegal actions by the action mask. The flags give the policy the information the masks use (Part~IV.7); the environment --- not the network --- remains the single authority on legality.

% ---------------------------------------------------------------------------
\subsection{M8 --- The policy and the two continuous heads}
\textbf{What.} \code{PhyloNetworkModelOneStep} keeps PhyloGFN's encoder and pair head and adds \code{single\_logits\_head}; \code{MergeEdgeModel} serves root-mode and two-edge-mode merges in one batch; \code{ReticulationModel} samples $(\log b_{hx},\mathrm{logit}\,\theta)$.
\begin{lstlisting}
pairs   = reps[:, rows] + reps[:, cols]                    # e_i + e_j   (PhyloGFN)
singles = reps                                             # e_i         (new)
pair_logits   = self.pair_logits_head(cat[pairs, summary]).squeeze(-1)
single_logits = self.single_logits_head(cat[singles, summary]).squeeze(-1)
logits = torch.cat([pair_logits, single_logits], dim=1).masked_fill(~action_mask, -inf)
actions = self.sample(logits, action_mask, random_spec)   # epsilon-greedy over LEGAL actions
log_paths_pf = log_softmax(logits)[arange(B), actions]
\end{lstlisting}
\begin{lstlisting}
class ReticulationModel(nn.Module):        # 2-D diagonal Gaussian mixture over (log b_hx, logit theta)
    def forward(self, summary_reps, lineage_reps, input_ret_actions=None):
        dist = self.distribution(torch.cat([summary_reps, lineage_reps], 1))
        z = dist.sample();  z = [z0.clip(-10, 0), z1.clip(-r-2, r+2)]
        lp = dist.log_prob(z)
        b_hx, theta = exp(z0), sigmoid(z1)
        lp = lp - z0 - log(theta) - log(1 - theta)          # change of variables for b and theta
        return b_hx, theta, lp
\end{lstlisting}
\textbf{Why.} The discrete choice and the continuous draw are two factors of one transition density; TB needs $\log\PF$ of the \emph{whole} step, so the head's log-density (with the Jacobian terms) is added to the discrete log-probability in the generator (M10). Squashing the mean of $\log b$ into $[-10,0]$ is PhyloGFN's trick (App.~F); the analogous squash of $\mathrm{logit}\,\theta$ into $[-6,6]$ keeps $\theta$ away from the numerically degenerate ends $0$ and $1$ at initialisation.

% ---------------------------------------------------------------------------
\subsection{M9 --- Exact class sizes and the topology prior}
\textbf{What.} \file{level1\_counts.py} implements the closed formula of Bouvel, Gambette \& Mansouri (2020, Prop.~5.4) for $r(n,k)$, the number of rooted binary level-1 networks with $n$ labelled leaves and $k$ reticulations, and the prior normaliser $\log Z_\lambda=\log\sum_{k\le R_{\max}}r(n,k)e^{-\lambda_Rk}$ (\code{EXP\_R}) or $\log\sum_{k\le R_{\max}}e^{-\lambda_Rk}$ (\code{UNIFORM\_R}); the environment also uses $\log r(N,R)$ in the reward under \code{UNIFORM\_R}.
\begin{lstlisting}
def log_topology_prior_normalizer(n, r_max, lambda_r):
    terms = [log(level1_count(n, k)) - lambda_r * k for k in range(0, r_max + 1)]
    mx = max(terms);  return mx + log(sum(exp(t - mx) for t in terms))
\end{lstlisting}
\textbf{Why.} It is the network analogue of PhyloGFN's \code{tree\_factor} $=-\log(2N-5)!!$ in the MLL estimator: without it the ``MLL'' would be off by the log of the prior normaliser. The formula was checked against the paper's Table~2 totals (36, 723, 20\,280, 730\,755) and against our own enumeration for $N\le6$; exact integers, no floating point.

% ---------------------------------------------------------------------------
\subsection{M10 --- The generator}
\textbf{What.} \code{TBGFlowNetNet.forward(input\_dict)} runs the policy, then routes the sampled ids: merges go to \code{MergeEdgeModel}, reticulations to \code{ReticulationModel}; the log-densities are added with \code{index\_add}; the environment-format action dicts are assembled. Replay trajectories arrive as \code{input\_actions} (one dict or \code{None} per trajectory) and are converted into forced ids and forced parameters (\code{nan} = not forced).
\begin{lstlisting}
is_merge = actions < env.n_pairs
sel = where(is_merge);  edges, lp = self.edges_model(summary[sel], left, right, last_step[sel], edge_given[sel])
log_pf = log_pf.index_add(0, sel, lp)
sel = where(~is_merge); b_hx, theta, lp = self.ret_model(summary[sel], reps[sel, lineage], ret_given[sel])
log_pf = log_pf.index_add(0, sel, lp)
# loss (unchanged form):  ( log Z + log_pf  -  log_reward - log_pb )^2
\end{lstlisting}
\textbf{Why.} PhyloGFN's generator always calls the edge model for \emph{every} row of the batch; ours must split rows by action type. The TB loss, $\log Z$ parametrisation (256-vector), optimiser groups, gradient clipping and temperature rescaling are copied unchanged.

% ---------------------------------------------------------------------------
\subsection{M11 --- Rollouts, replay buffer, evaluation}
\textbf{Rollout worker.}
\begin{lstlisting}
done = np.array([s.is_done for s in states])
while not done.all():
    active = np.where(~done)[0]
    input_dict = env.prepare_rollout_inputs(feats[active], act_states, random_spec, input_actions)
    ret = generator(input_dict)
    new_states, new_feats, ls, lr = env.batch_apply_actions(ret['actions'], feats[active], act_states)
    feats[active] = new_feats
    log_pf = log_pf.index_add(0, active, ret['log_paths_pf'])
    log_pb = log_pb.index_add(0, active, -log(tensor([s.num_parents() for s in new_states])))
    ... mark finished trajectories, store log_rewards / log_scores / n_ret
\end{lstlisting}
\textbf{Why.} Lock-step batching (\code{while tree\_features.shape[1] > 1}) is impossible with variable lengths; the \code{done} mask and per-trajectory accumulation are the minimal change. The parent count of the \emph{new} state is used, as in PhyloGFN; the last step now has one parent instead of $2N-3$ (rooted terminal object).\\
\textbf{Replay buffer} (\file{training\_data\_loader\_net.py}): same heap logic, keyed by the canonical network signature; backward trajectories from \code{sample\_backward\_from\_network}.\\
\textbf{Evaluator} (\file{gfn\_evaluator\_net.py}): eq.~\eqref{eq:mll} with $-\log Z_\lambda$ from M9, plus the histogram of sampled $R$ (the posterior over the number of reticulations) --- the quantity the validation plan calls ``posterior mode $R$'' (T8).

% ---------------------------------------------------------------------------
\subsection{M12 --- Registrations, configuration, scripts, tests}
\textbf{Modified original files (the entire diff to the original code).}
\begin{lstlisting}
# src/env/__init__.py
if cfg.ENV.ENVIRONMENT_TYPE == 'ONE_STEP_LEVEL1_NETWORK': env = PhyloNetworkEnv(cfg, all_seqs)
# src/gfn/build.py
if cfg.ENV.ENVIRONMENT_TYPE == 'ONE_STEP_LEVEL1_NETWORK': generator = TBGFlowNetNet(cfg.GFN, env, device, ddp)
# src/utils/utils.py  (correct_cfg_data): input size of the policy embedding
cfg.GFN.MODEL.TRANSFORMER.SEQ_EMB.INPUT_SIZE = 2 * seq_length * vocab_size + 4
# src/configs/defaults.py: ENV.NETWORK.{R_MAX, LAMBDA_R, PRIOR_TYPE}, TRANSFORMER.SINGLE_LOGITS_HEAD, RETICULATION_MODELING
\end{lstlisting}
\textbf{New scripts.} \pfile{train_net.py} (training loop without DDP/AMP; same schedules), \pfile{src/configs/network_cfgs/cfg_net_small_cpu.yaml} (a CPU-sized configuration), \pfile{tools/simulate_network_data.py} (alignments simulated on a known level-1 network; writes the truth), \pfile{tools/compare_tree_mode.py} ($R_{\max}=0$ versus PhyloGFN), \pfile{tests/test_network_env.py} (Part~VI).
\clearpage
